符号

.detach() 方法

@torch.inference\_mode() 装饰器

.unsqueeze() 操作

以编程方式加载 .jsonl 文件

@inference\_mode

--depth 1 选项

== 运算符

\textless{}think\textgreater{} 标签

.squeeze() 操作

@torch.no\_grad 装饰器

A

仅答案损失

自回归地

自回归文本生成

优势值跟踪

argmax 函数，第 2 处，第 3 处

优势值

通过……从采样轨迹中准备学习信号

显式添加

更多修改、提示与技巧，第 2 处

训练模型输出 \textless{}think\textgreater{} 词元，第 2 处

使用 \textless{}think\textgreater{} 词元，第 2 处

avg\_logprob\_score 函数

avg\_logprob\_answer 函数，第 2 处，第 3 处，第 4 处，第 5 处，第 6 处

B

基于基准测试的评估

brevity\_bonus，第 2 处

批量训练

build\_examples 函数，第 2 处

boxed\_bonus，第 2 处

BPE（字节对编码）

C

CoT（思维链）

checkpoint\_every

条件概率

compute\_grpo\_loss\_with\_kl 函数

计算图

传统 LLM（大语言模型）

compute\_grpo\_loss 函数，第 2 处，第 3 处，第 4 处，第 5 处，第 6 处，第 7 处，第 8 处

思维链提示，第 2 处

cross\_entropy 函数，第 2 处，第 3 处

检查点

保存的模型检查点，加载与评估，第 2 处

裁剪的策略比率

计算，第 2 处

通过……稳定序列级 GRPO，第 2 处

用……训练

count\_samples 函数，第 2 处，第 3 处

聊天机器人界面

compute\_example\_loss 辅助函数

封闭世界设定

compute\_length 函数

交叉熵损失

calc\_next\_token\_probas 函数，第 2 处

csv\_log\_path

calc\_next\_token\_logprobas 函数，第 2 处，第 3 处

compute\_grpo\_loss\_plus\_format\_reward 函数

D

蒸馏推理模型

评估蒸馏模型，第 2 处

为其加载 MATH 训练数据集，第 2 处

加载预训练模型

训练样本，第 2 处

草稿

数据类型

蒸馏

计算训练损失与验证损失，第 2 处

评估蒸馏模型，第 2 处

为其加载 MATH 训练数据集，第 2 处

加载预训练模型

后续步骤

在快速发展的领域中保持跟进

用于蒸馏的训练循环，第 2 处

数据集

为推理蒸馏而生成

datasets 库

蒸馏，第 2 处

概述，第 2 处

训练循环，第 2 处

面向高效推理的模型蒸馏

训练样本，第 2 处

download\_qwen3\_small() 函数

蒸馏损失

E

评估

加载评估数据集，第 2 处

评估

构建数学验证器

提取最终答案框，第 2 处

实现便于文本生成的封装器

加载评估数据集，第 2 处

加载预训练模型以生成文本，第 2 处

规范化提取的答案，第 2 处

验证数学等价性，第 2 处

evaluate\_json.py 脚本

evaluate\_math500\_stream 函数，第 2 处，第 3 处，第 4 处，第 5 处

评估准确率，计算

equality\_check 函数，第 2 处，第 3 处，第 4 处，第 5 处

提取的答案，规范化，第 2 处

熵跟踪，第 2 处

encode 方法

评估推理模型

构建数学验证器

提取最终答案框，第 2 处

实现便于文本生成的封装器

加载预训练模型以生成文本，第 2 处

ex 玩具示例

extract\_bonus，第 2 处

extract\_final\_candidate 函数，第 2 处，第 3 处

涌现特性

F

format\_distilled\_answer 函数

filter\_examples\_by\_max\_len 步骤

格式奖励

显式添加，第 2 处

完整前向传播

fulltext\_bonus

G

generate\_text\_stream\_concat 函数，第 2 处

generate\_text\_stream\_concat 模型

get\_last\_boxed 函数

GRPO（广义奖励预测目标）

generate\_text\_basic\_stream

grade\_answer 函数，第 2 处，第 3 处，第 4 处，第 5 处，第 6 处，第 7 处

generate\_stats 函数

get\_device() 函数

使用预训练 LLM 生成文本

编写最简文本生成函数，第 2 处

GRPO（广义响应策略优化）

跟踪性能指标，第 2 处，第 3 处，第 4 处

给答案评分，第 2 处

GRPO（广义带策略优化的 REINFORCE）流程

通过 GRPO 损失进行策略更新

GRPO（组相对策略优化）

改进

generate\_text\_basic\_stream 函数，第 2 处，第 3 处，第 4 处，第 5 处，第 6 处，第 7 处，第 8 处

generate\_text\_temp\_stream\_cache 函数

generate\_text\_temp\_stream\_cache

GRPO（组相对策略优化）

使用强化学习训练推理模型，第 2 处

generate\_text\_stream\_concat\_flex

generate\_text\_stream\_cache 函数

generate\_text\_stream\_concat\_flex 函数，第 2 处，第 3 处，第 4 处，第 5 处

generate\_text\_stream\_cache

generate\_text\_top\_p\_stream\_cache 函数

generate\_text\_basic\_stream\_cache

GRPO（组相对策略优化）算法

通过厨师类比理解 GRPO 的高层直觉

GRPO 高层流程

使用，第 2 处

gtruth\_answer 字段

GRPO（目标—奖励—提示优化）

用序列对数概率对采样轨迹评分，第 2 处

GRPO（组相对策略优化）训练循环

实现，第 2 处

grad\_clip\_norm

generate\_text\_basic\_stream\_cache 函数，第 2 处，第 3 处，第 4 处

使用……生成文本

编写最简文本生成函数，第 2 处

标准答案

贪心解码

GRPO（广义采样轨迹策略优化）

跟踪性能指标，第 2 处，第 3 处，第 4 处

H

硬蒸馏

启发式

Hugging Face 模型中心

硬件

需求与建议，第 2 处

heuristic\_score 函数，第 2 处，第 3 处

I

改进用于强化学习的 GRPO

显式添加格式奖励，第 2 处

更多修改、提示与技巧，第 2 处

训练模型输出 \textless{}think\textgreater{} 词元，第 2 处

使用 \textless{}think\textgreater{} 词元，第 2 处

指令微调

推理计算扩展

改进 GRPO 用于

显式添加格式奖励，第 2 处

用 KL 项控制模型变化幅度，第 2 处

更多修改、提示与技巧，第 2 处

使用裁剪的策略比率稳定序列级 GRPO，第 2 处

训练模型输出 \textless{}think\textgreater{} 词元，第 2 处

使用 \textless{}think\textgreater{} 词元，第 2 处

改进推理，使用

思维链提示，第 2 处

加载预训练模型，第 2 处

通过温度参数对词元分数（logits）重新缩放

推理

通过 KV 缓存加速，第 2 处

通过 PyTorch 模型编译加速，第 2 处

推理时扩展（inference scaling）

通过推理时扩展改进

为文本生成函数添加温度缩放

思维链提示，第 2 处

用温度缩放控制输出多样性

加载预训练模型，第 2 处

通过温度参数对词元分数（logits）重新缩放

从概率分布中采样下一个词元

选择 top-p 词元的子集

推理时扩展，第 2 处，第 3 处，第 4 处

为文本生成函数添加温度缩放

用温度缩放控制输出多样性

似然与概率

对数概率，第 2 处，第 3 处，第 4 处

下一个词元选择过程

从概率分布中采样下一个词元

自洽性采样，第 2 处

自我修正，第 2 处

自我修正循环，第 2 处

词元概率分数，第 2 处，第 3 处，第 4 处

top-k 过滤，第 2 处

top-p 过滤，第 2 处

改进用于强化学习

显式添加格式奖励，第 2 处

计算裁剪的策略比率，第 2 处

用 KL 项控制模型变化幅度，第 2 处

更多修改、提示与技巧，第 2 处

使用裁剪的策略比率稳定序列级 GRPO，第 2 处

训练模型输出 \textless{}think\textgreater{} 词元，第 2 处

使用裁剪的策略比率训练

使用 \textless{}think\textgreater{} 词元，第 2 处

iterations 参数

推理时计算扩展，第 2 处

推理时

为文本生成函数添加温度缩放

指令微调

J

联合概率

基于评判的评估

K

KL 损失项

KL（Kullback–Leibler）散度

用 KL 项控制模型变化幅度，第 2 处

实现 KL 损失项，第 2 处

使用 KL 损失项训练，第 2 处

kl\_loss

知识蒸馏

KV 缓存

通过……加速推理，第 2 处

KL 散度

KVCache 对象

kl\_coeff

L

load\_math\_train 函数

lr（学习率），第 2 处

log\_every

logits

通过温度参数对词元分数（logits）重新缩放

LLM（大语言模型）

用于文本生成

通过训练与推理技术改进推理，第 2 处

为其准备输入文本，第 2 处

使用基于规则的分数为回复打分，第 2 处

训练流程，第 2 处

对数概率评分（logprob scoring）方法

对数概率，第 2 处

用……评估模型置信度，第 2 处

用序列对数概率对采样轨迹评分，第 2 处

logprob

Latex 类，第 2 处

对数概率（logprobs）

load\_model\_and\_tokenizer

load\_state\_dict 方法

似然，第 2 处

load\_math500\_test 函数

LLM 作为评判者

学习信号

通过优势值从采样轨迹中准备

load\_model\_and\_tokenizer 函数

M

MATH（数学）数据集

为其加载 MATH 训练数据集，第 2 处

message\_thinking 字段

mini\_eval\_demo 函数，第 2 处

mini\_eval\_demo(model, tokenizer, device) 函数

多奖励训练

数学等价性，验证，第 2 处

message\_content 字段

最大化

中期训练阶段

Math 类

max 函数

最小化

数学验证器

模型蒸馏

概述，第 2 处

max\_new\_tokens

N

next\_ids 张量

规范化提取的答案，第 2 处

核采样

num\_rollouts，第 2 处

n\_kept 变量

next\_token\_probas

下一个词元选择过程

下一个词元概率

next\_token\_logits 张量，第 2 处

normalize\_text 函数，第 2 处，第 3 处

O

开放世界设定

概述

RLHF

RLVR

P

模式匹配

逻辑推理与之对比

PyTorch

通过模型编译进行推理，第 2 处

概率分数

plt.bar 柱状图

pprint 库

print 函数

plt.plot 折线图

策略更新

通过 GRPO 损失

PPO（近端策略优化）

逐词元概率

预训练阶段

过程奖励

prompt 参数

pprint 标准库

策略梯度损失，数学定义

提示模板

偏好微调，第 2 处

预训练模型，加载，第 2 处

概率

似然与之对比

策略优化算法

plot\_grpo\_metrics 函数

预训练模型，加载，第 2 处

pip 包安装器，第 2 处

pred\_answer 最终候选答案

《近端策略优化算法》

plot\_scores\_bar 函数

pyproject.toml 文件

partial 函数

后训练阶段

pprint 函数

Q

Qwen3Model 类

R

推理模型

推理训练

raw\_prompt

re 库

使用 RLVR 的推理模型

加载并评估保存的模型检查点，第 2 处

RLHF（基于人类反馈的强化学习），第 2 处，第 3 处

RL（强化学习），第 2 处

概述，第 2 处

reward\_format 函数

正则表达式

采样轨迹

通过优势值从中准备学习信号

采样，第 2 处

用序列对数概率评分，第 2 处

run\_demos\_table 函数

requirements.txt 文件，第 2 处

run\_demos\_table() 函数

奖励

计算

reasoning-from-scratch Python 包

reasoning-from-scratch 库

推理

通过推理时扩展改进

reward\_avg

推理蒸馏

为其生成数据集

基于规则的评分函数

推理模型，第 2 处

LLM 与人类推理

从零构建

定义，第 2 处

评估，第 2 处

未来方向

给答案评分，第 2 处

通过训练与推理技术改进推理，第 2 处

逻辑推理与之对比

模式匹配与逻辑推理的比较

自洽性采样，第 2 处

在没有显式规则的情况下模拟推理

使用强化学习训练，第 2 处，第 3 处

强化学习

改进 GRPO

用……训练推理模型，第 2 处，第 3 处

run\_demos\_table(tests) 函数

regex（正则表达式）

采样轨迹采样

构建推理模型的路线图

reward\_rlvr 函数，第 2 处，第 3 处

使用强化学习的推理模型

通过 GRPO 损失进行策略更新

在单个 GRPO 函数中整合所有内容，第 2 处

results 字典，第 2 处，第 3 处

RLVR（基于可验证奖励的强化学习），第 2 处

使用 GRPO，第 2 处

RLVR（带验证与重训练的强化学习）

加载并评估保存的模型检查点，第 2 处

results\_logprob 字典

原始字符串

S

序列对数概率

用……为采样轨迹评分，第 2 处

自洽性，第 2 处

SumBackward0 条目

缩放

通过温度参数对词元分数（logits）缩放

评分

使用基于规则的分数为 LLM 回复打分，第 2 处

为模型回复评分并迭代改进

序列似然

软蒸馏

sequence\_logprob\_and\_entropy 函数，第 2 处，第 3 处

softmax 函数

LLM 顺序文本生成过程，第 2 处

采样轨迹采样，第 2 处

统计推断

sequence\_logprob 函数，第 2 处，第 3 处，第 4 处，第 5 处

split\_into\_parts 函数

监督微调（SFT）

score\_fn 选项

self\_consistency\_vote 函数，第 2 处，第 3 处，第 4 处，第 5 处，第 6 处，第 7 处

SymPy 数学库

score\_fn 调用

self\_refinement\_loop 函数，第 2 处

sympy\_parser 函数，第 2 处

压缩张量

为采样轨迹评分

用序列对数概率，第 2 处

SFT（监督微调），第 2 处

评分器

自我修正，第 2 处，第 3 处，第 4 处，第 5 处

加载预训练模型，第 2 处

循环，第 2 处

使用基于规则的分数为 LLM 回复打分，第 2 处

为模型回复评分并迭代改进，第 2 处

用对数概率评估模型置信度，第 2 处，第 3 处，第 4 处

scale\_logits\_by\_temperature 函数

自洽性采样，第 2 处

sample\_response 函数，第 2 处，第 3 处，第 4 处

T

生成文本，使用

通过 KV 缓存加速推理，第 2 处

通过 PyTorch 模型编译加速推理，第 2 处

通过 PyTorch 模型编译进行推理，第 2 处

加载预训练模型，第 2 处

搭建编程环境，第 2 处

分词

数据集

top-p 设置，第 2 处

词元级似然

torch.max() 函数

词元概率分数，第 2 处

使用 RL 训练

RLHF

RLVR

概述，第 2 处

torch.argmax 函数

tokenizers 库

torch.argmax() 函数

torch.log\_softmax 函数，第 2 处

top-p 采样器

测试时计算扩展

训练子集

加载，第 2 处

分词器

top-p 采样

torch.cumsum 函数

torch.softmax() 函数

目标序列

训练样本，第 2 处

过滤数据集，第 2 处，第 3 处，第 4 处，第 5 处，第 6 处

格式化数据集，第 2 处，第 3 处

划分数据集，第 2 处，第 3 处，第 4 处，第 5 处，第 6 处

分词器，第 2 处，第 3 处

对数据集分词，第 2 处，第 3 处

使用强化学习训练

计算奖励

实现 GRPO 训练循环，第 2 处

加载 MATH 训练子集，第 2 处

加载预训练模型

通过 GRPO 损失进行策略更新

在单个 GRPO 函数中整合所有内容，第 2 处

采样轨迹采样，第 2 处

训练

使用强化学习的推理模型，第 2 处

用序列对数概率对采样轨迹评分，第 2 处

torch.compile

top-p 范围

top\_p\_filter 函数，第 2 处

温度参数

通过它重新缩放词元分数（logits）

温度设置，第 2 处，第 3 处

tokenizer-base.json 文件

top\_p 设置

tokenizer.eos\_token\_id 属性

top-k 过滤

训练循环

用于蒸馏，第 2 处

词元概率

文本生成

用于……的 LLM

为其添加温度缩放

硬件需求与建议，第 2 处

加载预训练模型，第 2 处

为 LLM 准备输入文本，第 2 处

LLM 顺序文本生成过程，第 2 处，第 3 处，第 4 处

搭建编程环境，第 2 处

训练推理模型

计算奖励

通过优势值从采样轨迹中准备

用序列对数概率对采样轨迹评分，第 2 处

torch.softmax 函数，第 2 处，第 3 处

torch.log 函数

跟踪性能指标

优势值跟踪

计算评估准确率，第 2 处

熵跟踪，第 2 处

执行训练运行，第 2 处

检查训练运行，第 2 处，第 3 处，第 4 处

绘制其他指标

train\_distillation 函数

train\_rlvr\_grpo() 函数

训练损失

计算，第 2 处

分词器，第 2 处

词元对数概率

温度缩放，第 2 处

用……训练推理模型

计算奖励

实现 GRPO 训练循环，第 2 处

加载 MATH 训练子集，第 2 处

加载预训练模型

通过 GRPO 损失进行策略更新

通过优势值从采样轨迹中准备

在单个 GRPO 函数中整合所有内容，第 2 处

采样轨迹采样，第 2 处

用序列对数概率对采样轨迹评分，第 2 处

torch.exp() 函数

top-p 过滤器

训练轮次，第 2 处

torch.log\_softmax() 函数，第 2 处，第 3 处

t\_idx 索引张量

top-p 过滤，第 2 处

torch.multinomial 函数，第 2 处

top-p 词元

选择其子集

torch.prod 函数

torch.multinomial() 函数

测试时扩展

使用 RLVR 训练

加载并评估保存的模型检查点，第 2 处

温度设置

使用强化学习训练推理模型

在单个 GRPO 函数中整合所有内容，第 2 处

词元，第 2 处

\textless{}think\textgreater{} 词元，第 2 处

torch 库

使用 KL 损失项训练

KL 方向与采样轨迹采样的对比

用 KL 项提升训练稳定性

U

解压张量

使用 GRPO

通过厨师类比理解 GRPO 的高层直觉

GRPO 高层流程

uv Python 包

V

验证损失

可验证奖励

验证损失

计算，第 2 处

验证数学等价性，第 2 处

W

使用强化学习

加载 MATH 训练子集，第 2 处

采样轨迹采样，第 2 处

使用预训练 LLM

通过 KV 缓存加速推理，第 2 处

通过 PyTorch 模型编译加速推理，第 2 处
